\documentclass{article}
\usepackage{amsmath, amsthm, amssymb}
\usepackage{hyperref}
\usepackage[margin=1in]{geometry}

\newtheorem{theorem}{Theorem}
\newtheorem{proposition}[theorem]{Proposition}
\newtheorem{lemma}[theorem]{Lemma}
\newtheorem{corollary}[theorem]{Corollary}
\theoremstyle{definition}
\newtheorem{definition}[theorem]{Definition}
\newtheorem{observation}[theorem]{Observation}
\newtheorem{example}[theorem]{Example}
\theoremstyle{remark}
\newtheorem{remark}[theorem]{Remark}

\title{The Bi Cluster Test at $n=4$: Partial No, Structural Yes}
\author{Clio}
\date{2026-04-21}

\begin{document}
\maketitle

\begin{abstract}
We test whether the principal minors $d_k(n) = \det M^{(k)}[D_n]$ of the staircase Hecke element $\Pi_q$ sit inside Bi's cluster algebra on the Hernandez--Leclerc subcategory $\mathcal{C}_{v,\beta}$. At $n=4$ the direct identification $d_k \leftrightarrow D_k$ fails on three counts: count, exchange form, and ambient ring. The failure is diagnostic: the staircase case $v = \delta(\beta)$ is precisely the degenerate case of Bi's framework, where the new content collapses to Gei\ss--Leclerc--Schr\"oer. Structural matches in rank and generator count suggest the live alternative is Sandvik's monodromic-Hecke / Soergel-bimodule route.
\end{abstract}

\section{Setup --- the question}

Bi (\href{https://arxiv.org/abs/2602.11559}{arXiv:2602.11559}) constructs, for each reduced word $\beta$ of $w_0 \in W$ and each $v \leq \delta(\beta)$ in the higher Bruhat order, a monoidal subcategory $\mathcal{C}_{v,\beta}$ of the Hernandez--Leclerc category whose Grothendieck ring $K_0(\mathcal{C}_{v,\beta})$ carries a cluster algebra structure. The mutable cluster variables are (quantum) minors of Bott--Samelson type; frozen variables are boundary Pl\"ucker coordinates.

Our natural test case: the staircase reduced word
$$w_0 \;=\; s_1 s_2 s_3 \cdot s_1 s_2 \cdot s_1 \qquad \text{in } S_4 \;=\; W(A_3),$$
with length $\ell(w_0) = 6 = |D_4|$ (the descent-paired lattice size at $n=4$). The \emph{staircase Hecke element} is
$$\Pi_q \;=\; (1+T_1)(1+T_2)(1+T_1)(1+T_3)(1+T_2)(1+T_1) \;\in\; H_q(S_4).$$
It is the Bott--Samelson element at this reduced word. Its principal minors on the descent-paired set $D_n \subset S_n$ have been the central object of our April 2026 work; the determinantal identity
$$d_{2j+1}(n) \cdot d_{2j-1}(n) \;=\; q^{|D_n|} \, d_{2j}(n)^2 \qquad \text{(Hirota-type)}$$
has been proved at $j=1$ and $j=2$ and heavily supported at $j=3$.

We asked two sharp questions:

\begin{quote}
\textbf{H1.} Are the $d_k(n)$ individual cluster variables in Bi's seed on the staircase?

\smallskip
\textbf{H2.} Is the Hirota identity $d_{2j+1} d_{2j-1} = q^{|D_n|} d_{2j}^2$ an exchange relation?
\end{quote}

A ``yes'' would identify $\Pi_q$'s integrable structure with the KKKO/Bi cluster structure on $K_0(\mathcal{C}_{v,\beta})$; it would be a categorical origin for the palindromic factor.

\section{Data at $n=4$}

Direct Hecke computation (see Appendix A) gives
\begin{align}
d_1(4) &\;=\; q^6, \\
d_2(4) &\;=\; q^9 (q+1)^3 (2q+1), \\
d_3(4) &\;=\; q^{18} (q+1)^6 (2q+1)^2 \;=\; d_2(4)^2.
\end{align}
The first $q$-shifted pair identity checks:
$$d_3(4) \cdot d_1(4) \;=\; q^{18}(q+1)^6(2q+1)^2 \cdot q^6 \;=\; q^6 \cdot q^{18}(q+1)^6(2q+1)^2 \;=\; q^6 \cdot d_2(4)^2,$$
with shift exponent $|D_4| = 6$, consistent with the general Hirota conjecture.

So at $n=4$ there are only \emph{three} distinct $d_k$ quantities, and the middle one has a clean factorisation into three ``staircase atoms'' $q$, $(q+1)$, $(2q+1)$.

\section{Bi's initial seed for the staircase $w_0 = (1,2,3,1,2,1)$ in $A_3$}

From Bi \S 2 and Theorem 3.1, the initial seed for $\mathcal{C}_{w_0,\beta}$ at the staircase word has six vertices, with the standard frozen choice $\{3,5,6\}$ (boundary Pl\"uckers) and mutable $\{1,2,4\}$. The cluster variables are the six positive-root generalised minors of $\mathbb{C}[N_-]$ for $\mathrm{SL}_4$, matching Gei\ss--Leclerc--Schr\"oer:
\begin{align}
D_1 &\;=\; \Delta_{2,1} \\
D_2 &\;=\; \Delta_{23,12} \\
D_3 &\;=\; \Delta_{234,123} \qquad \text{(frozen, size 3)} \\
D_4 &\;=\; \Delta_{3,2} \\
D_5 &\;=\; \Delta_{34,23} \qquad \text{(frozen, size 2)} \\
D_6 &\;=\; \Delta_{4,3} \qquad \text{(frozen, size 1)}
\end{align}
By Pl\"ucker size: three size-1 minors, two size-2 minors, one size-3 minor. This is the $A_3$ positive-root system organised by height.

The exchange relation (Bi eq.~4.16) is the standard three-term Fomin--Zelevinsky form:
$$D_{k-1} \cdot \mu_{k-1}(D_{k-1}) \;=\; q^A \, D_k \cdot D_{(k-1)^-} \;+\; q^B \, D_{k+1} \cdot D_{k^-}.$$

\begin{remark}
For $v = w_0 = \delta(\beta)$ --- the full staircase --- Bi's new content collapses back to the classical GLS seed on $\mathbb{C}[N_-]$ (Theorem 3.11, final sentence). In Bi's framework, the staircase $v = w_0$ is the \emph{degenerate / boundary} case. Bi's innovation operates for $v < \delta(\beta)$, where the seed is genuinely smaller than the full GLS seed and the cluster structure detects sub-Bruhat-order information.
\end{remark}

\section{Three mismatches}

\begin{observation}[Count]
Bi's seed at $v = w_0$ in $A_3$ has \textbf{6} cluster variables. We have only \textbf{3} distinct $d_k$ quantities at $n=4$. If $d_k$ were each an individual cluster variable, we would need three variables out of six. No natural assignment presents itself: the $d_k$ have no canonical subset structure inside the positive-root system.
\end{observation}

\begin{observation}[Exchange form]
Bi's exchange is three-term (two monomials on the right). Our Hirota identity is two-term bilinear:
$$d_{2j+1} \, d_{2j-1} \;=\; q^{|D_n|} \, d_{2j}^2.$$
A two-term bilinear identity cannot arise from a single mutation of a three-term exchange unless one of Bi's monomials vanishes identically --- which does not happen in a generic cluster algebra. Even allowing aggregation into products of several $D_k$, the right-hand side of a cluster exchange never collapses to a single monomial in the generic case.
\end{observation}

\begin{observation}[Ambient ring]
Bi's $D_k$ live in the quantum unipotent coordinate ring $\mathbb{C}_q[N_-]$, a quotient of $U_q(\mathfrak{n}_-)$. Our $d_k$ live in the Hecke algebra $H_q(S_n)$. These are related at the dual level by Schur--Weyl duality
$$H_q(S_n) \curvearrowright V^{\otimes n} \curvearrowleft U_q(\mathfrak{sl}_n),$$
but the pairing between Bott--Samelson Hecke elements and Pl\"ucker-type minors is nontrivial. The natural image of a Bott--Samelson element of $H_q$ in $\mathbb{C}_q[N_-]$ --- e.g.\ via the Kashiwara--Lusztig global basis --- is typically a \emph{product} of several generalised minors, not a single one.
\end{observation}

These three observations together kill H1 and H2 in the literal reading $d_k \leftrightarrow D_k$.

\section{Structural positives (to guide next experiment)}

The failure has a pattern, and the pattern is informative.

\begin{observation}[Rank-count match]
$|D_4| = 6$ equals the number of Bi cluster variables, which equals $\ell(w_0) = |\Phi^+(A_3)|$. The descent-paired lattice size equals the Bi seed size equals the number of positive roots. Both sides count the generators of the Bott--Samelson / positive-root system, so this is not coincidence --- but it is the \emph{outer} count matching, not the individual generators.
\end{observation}

\begin{observation}[Size distribution]
Bi's six generators split by Pl\"ucker size as $(3, 2, 1)$: three size-1, two size-2, one size-3. Our $d_k$ split by level as $(k=1, k=2, k=3)$, i.e.\ three levels with one quantity each. Both are ``staircase'' sequences, but with opposite monotonicity (Bi's is decreasing, ours is uniform-by-level). The natural question: is there a bijection pairing $d_k$ with the aggregated monomial over size-$(n-k)$ minors? Numerically this would mean
$$d_k(4) \;\stackrel{?}{=}\; \prod_{\text{size-}(n-k) \text{ minors } D_i} D_i^{a_i}$$
for some exponents $a_i$, evaluated at the GLS parametrisation. This is testable but was not carried out here.
\end{observation}

\begin{observation}[Bott--Samelson interpretation]
$\Pi_q$ is precisely the Bott--Samelson element of $H_q$ at the reduced word $(1,2,1,3,2,1)$. In the Soergel-bimodule category this is $BS(w_0)$, whose class in $K_0(\text{SBim}) \cong H_q$ is the product of Kazhdan--Lusztig generators $C'_{s_i}$. This is the \textbf{Sandvik / monodromic Hecke} route (\href{https://arxiv.org/abs/2604.15582}{arXiv:2604.15582}), not the Bi cluster route.

The two categorical survivors from our three-routes analysis operate at different levels of the geometric Langlands tower:
\begin{itemize}
\item \textbf{Bi:} quantum-group reps, $K_0$ of a monoidal subcategory of finite-dim $U_q(\widehat{\mathfrak{sl}_n})$-modules.
\item \textbf{Sandvik:} Soergel bimodules / parity sheaves on $G/B$, $K_0$ of Hecke-module categories.
\end{itemize}
$\Pi_q$ is a Hecke-algebra element \emph{on the nose}, so the Soergel / parity-sheaf side is the direct home.
\end{observation}

\section{Conclusion}

Both H1 and H2 are \textbf{false} at $n=4$ under the direct identification $d_k \leftrightarrow D_k$. The numerical count (6 vs 3), the exchange form (three-term vs two-term), and the ambient ring ($\mathbb{C}_q[N_-]$ vs $H_q(S_n)$) each independently rule out the naive correspondence.

The reason is structural: the staircase case $v = \delta(\beta)$ is precisely the \emph{degenerate} case of Bi's framework, where the seed reduces to the classical GLS seed and Bi's new content --- the $v < \delta(\beta)$ refinement --- is absent. We chose the test case that sits at the boundary of Bi's theorem, which is exactly where it has least to say.

The negative answer is diagnostic. Candidates remaining:
\begin{enumerate}
\item[(a)] $d_k$ as a \emph{product / aggregation} of Bi cluster variables at a fixed Pl\"ucker size --- requires computing Bi classes in $\mathbb{Z}[q]$ and matching the $(q+1)^3(2q+1)$ factorisation of $d_2(4)$ against products of GLS minors.
\item[(b)] $d_k$ as a Grothendieck class in a different monoidal category --- \textbf{Sandvik monodromic Hecke}, where $BS(w_0)$ is a parity sheaf and self-duality gives the palindromic factor a priori. This is the structurally natural home.
\item[(c)] The Hirota identity as the \emph{shadow} of a higher-level identity in Bi's framework (double mutation, interval, product seed) at $v < \delta(\beta)$, projected back to the staircase by specialisation. More speculative.
\end{enumerate}

The cleanest next experiment is (b): compute Soergel-bimodule classes for $BS(w_0)$ and check whether the Hirota identity is a shadow of a Grothendieck-group relation in $K_0(\text{SBim})$. This is consistent with the \emph{KL basis unification} theme of the three April 2026 papers, and it is the route where $\Pi_q$ appears on the nose rather than after a Schur--Weyl transposition.

\appendix
\section{Scripts and data}

\begin{itemize}
\item \verb|/home/clio/projects/scratch/2026-04-21-staircase-n4-bi-data.py| --- direct Hecke computation of $d_1(4), d_2(4), d_3(4)$ and the first-pair identity check.
\item \verb|/home/clio/projects/scratch/2026-04-21-staircase-n4-data.md| --- tabulation of the staircase data at $n=4$ in human-readable form.
\item \verb|/home/clio/projects/scratch/2026-04-21-bi-cluster-extract.md| --- extract of Bi \S 2--3 with the initial-seed specialisation to $v = w_0$ in $A_3$.
\end{itemize}

\end{document}
